\documentclass[11pt,a4paper]{article}

\usepackage[margin=24mm]{geometry}
\usepackage{kotex}
\usepackage{amsmath,amssymb,mathtools,bm}
\usepackage{booktabs,longtable,array}
\usepackage[hidelinks]{hyperref}
\usepackage{setspace}

\setstretch{1.14}
\setlength{\parindent}{0pt}
\setlength{\parskip}{0.45em}

\newcommand{\dd}{\mathrm{d}}
\newcommand{\ext}{\mathrm{ext}}
\newcommand{\cell}{\mathrm{cell}}
\newcommand{\bg}{\mathrm{bg}}
\newcommand{\tot}{\mathrm{tot}}
\newcommand{\eqs}{\mathrm{eq}}
\newcommand{\app}{\mathrm{app}}
\newcommand{\obs}{\mathrm{obs}}
\newcommand{\OCV}{\mathrm{OCV}}
\newcommand{\absI}{I_{\mathrm{abs}}}
\newcommand{\sgn}{\mathrm{sgn}}

\hypersetup{
  pdftitle={Charge-Balance Theory for Graphite ICA dQdV Peak Shape},
  pdfauthor={},
  pdfsubject={Graphite ICA dQdV theoretical derivation}
}

\begin{document}

\begin{center}
{\LARGE\bfseries 전하 보존식 기반 흑연 ICA/dQdV 이론 전개}\\[0.6em]
{\large 상변이 피크 면적, 꼬리 형상, 온도 의존성의 논리적 연결}
\end{center}

\begin{abstract}
이 문서는 리튬이온전지 흑연 음극의 incremental capacity analysis, 즉 ICA \(=\dd Q/\dd V\)에서 상변이 피크가 왜 나타나고, 같은 상변이의 총 면적은 보존되면서도 온도와 현재 전위 상태에 따라 꼬리 형상이 왜 달라지는지를 전하 보존식에서 출발해 전개한다. 핵심은 전압 곡선을 먼저 가정하지 않는 것이다. 외부에서 흘린 전하 \(Q_{\ext}\), 정규화 전하 좌표 \(q\), 흑연 내부 상변이 진행률 \(\xi_j\), 배경 저장량 \(Q_{\bg}\)를 먼저 정의하고, 내부 음극 전위 \(V_n\)는 이들이 동시에 만족해야 하는 전하 보존식의 해로 정의한다. 그 뒤에 평형 OCV, 비평형 상변이 지연, \(dV/\dd q\), \(dQ/\dd V\), 피크 면적과 꼬리 형상을 차례로 유도한다. 따라서 피크 면적은 전이 용량과 진행률 변화량으로 정해지고, 꼬리 길이는 평형 위치 자체가 아니라 진행률이 평형값을 따라가는 속도, 즉 온도 의존적인 relaxation 길이로 설명된다.
\end{abstract}

\tableofcontents
\newpage

\section{설명해야 할 현상}

흑연 음극 ICA에서 관측되는 피크는 임의로 그린 함수가 아니라 흑연 내부의 상변이 또는 staging transition이 전압에 따라 진행되면서 나타나는 미분 신호이다. 사용자가 확인한 핵심 현상은 다음과 같이 정리할 수 있다.

\begin{enumerate}
\item 같은 상변이가 완전히 진행된다면 그 피크가 담당하는 총 전하량, 즉 피크 면적은 같은 값에 가까워야 한다.
\item 그러나 피크의 꼬리 부분은 온도와 현재 전위 상태에 따라 달라진다.
\item 낮은 온도에서는 꼬리가 길게 늘어지고, 높은 온도에서는 꼬리가 상대적으로 빨리 끝난다.
\end{enumerate}

이 세 문장을 동시에 만족하려면 ``피크의 면적''과 ``피크의 모양''을 분리해야 한다. 면적은 얼마만큼의 상변이 용량이 지나갔는지를 말한다. 모양은 그 용량이 전압축 위에 어떻게 분포했는지를 말한다. 수학적으로는 같은 적분값을 갖는 함수라도 좁은 구간에 집중되면 높고 날카로운 피크가 되고, 넓은 구간에 분포하면 낮고 긴 꼬리를 갖는다. 실제 수식으로는
\[
\text{면적}=\int \frac{\dd Q}{\dd V}\,\dd V,
\qquad
\text{모양}=\frac{\dd Q}{\dd V}\text{가 }V\text{에 따라 분포하는 방식}
\]
이다.

따라서 이 문서의 목적은 피크 함수를 먼저 고르는 것이 아니다. 먼저 흑연 내부의 전하 저장과 상변이 진행률을 정의하고, 그 결과로 \(dQ/\dd V\) 형상이 어떻게 나오는지 보이는 것이다.

\section{측정 좌표와 ICA의 정의}

분석 구간을 \(t_0\le t\le t_1\)라고 하자. 전류의 부호는 충전과 방전 convention에 따라 달라지므로, 이 문서에서는 선택한 분석 방향으로 흘러간 전류 크기를
\begin{equation}
\absI(t)=|I(t)|
\label{eq:current_magnitude}
\end{equation}
로 둔다. 외부 회로를 통해 분석 시작점부터 시간 \(t\)까지 흘린 전하량은
\begin{equation}
Q_{\ext}(t)=\int_{t_0}^{t}\absI(t')\,\dd t'
\label{eq:qext_definition}
\end{equation}
이다. 기준 용량 \(Q_{\cell}>0\)로 나누면 무차원 진행 좌표
\begin{equation}
q(t)=\frac{Q_{\ext}(t)}{Q_{\cell}}
\label{eq:q_definition}
\end{equation}
를 얻는다. \(Q_{\cell}\)은 수식에서 쿨롱 단위로 사용한다. 만약 실험값이 Ah 단위이면
\begin{equation}
Q_{\cell}[\mathrm{C}]=3600\,Q_{\cell}[\mathrm{Ah}]
\label{eq:ah_to_c}
\end{equation}
로 바꾸어야 한다.

\(\absI>0\)인 구간에서는 식 \eqref{eq:qext_definition}과 \eqref{eq:q_definition}으로부터
\begin{equation}
\frac{\dd q}{\dd t}=\frac{\absI}{Q_{\cell}},
\qquad
\frac{\dd}{\dd q}=\frac{Q_{\cell}}{\absI}\frac{\dd}{\dd t}
\label{eq:time_to_q}
\end{equation}
가 된다. 휴지 구간에서는 \(\absI=0\)이므로 \(q\)가 시간에 따라 변하지 않는다. 따라서 휴지 relaxation은 \(q\) 미분식이 아니라 시간식으로 다루어야 한다.

측정 전압 또는 분석에 쓰는 apparent 전압을 \(V_{\obs}\)라고 하자. ICA와 DVA는 먼저 측정량 수준에서 다음처럼 정의된다.
\begin{equation}
\mathrm{ICA}(V_{\obs})
=
\frac{\dd Q_{\ext}}{\dd V_{\obs}},
\qquad
\mathrm{DVA}(Q_{\ext})
=
\frac{\dd V_{\obs}}{\dd Q_{\ext}}.
\label{eq:ica_dva_measurement}
\end{equation}
이 정의는 아직 흑연 상변이를 가정하지 않는다. 단지 ``전압이 조금 변할 때 외부 전하가 얼마나 변했는가''와 그 역수를 말할 뿐이다.

\section{흑연 내부 저장량의 분해}

흑연 음극의 전하 저장은 두 부분으로 나누어 쓴다. 첫째는 뚜렷한 피크로 분리하지 않는 배경 저장량이고, 둘째는 분리 가능한 effective transition들의 저장량이다.

배경 저장량을
\begin{equation}
Q_{\bg}=Q_{\bg}(V,T)
\label{eq:qbg_definition}
\end{equation}
라고 둔다. 여기서 \(V\)는 아직 확정된 전위가 아니라 임시 전압 변수이다. \(Q_{\bg}\)는 그래프를 보기 좋게 만들기 위한 baseline이 아니라 실제로 전하 보존식에 들어가는 저장량이다.

흑연 내부의 effective transition을 \(j=1,\ldots,N_p\)로 표시한다. 각 transition의 진행률은
\begin{equation}
0\le \xi_j\le 1
\label{eq:xi_bounds}
\end{equation}
이다. \(\xi_j=0\)은 해당 transition이 분석 방향으로 아직 진행되지 않은 상태이고, \(\xi_j=1\)은 완료된 상태이다. 해당 transition이 담당할 수 있는 총 전하량을 \(Q_{j,\tot}>0\)라고 쓰면, transition 저장량은
\begin{equation}
Q_{\mathrm{tr}}=\sum_{j=1}^{N_p}Q_{j,\tot}\xi_j
\label{eq:transition_storage}
\end{equation}
이다.

이제 내부 상태 벡터를
\begin{equation}
\bm{\xi}=(\xi_1,\ldots,\xi_{N_p})
\label{eq:xi_vector_ch1}
\end{equation}
로 둔다. 중요한 점은 \(q\)와 \(\xi_j\)가 같은 것이 아니라는 점이다. \(q\)는 외부 회로를 지난 전하의 좌표이고, \(\xi_j\)는 흑연 내부의 특정 transition이 얼마나 진행되었는지를 나타내는 상태 변수이다. 같은 \(q\)에서도 온도, 전류, 이전 이력에 따라 \(\bm{\xi}\)가 달라질 수 있다.

\section{전하 보존식이 전압보다 먼저 온다}

외부에서 흘린 전하량은 흑연 내부 어딘가에 저장되어야 한다. 따라서 임시 전압 \(V\)에 대해 전하 보존 residual을
\begin{equation}
\mathcal{G}(V;q,T,\bm{\xi})
=
Q_{\bg}(V,T)
+\sum_{j=1}^{N_p}Q_{j,\tot}\xi_j
-Q_{\cell}q
\label{eq:residual_definition}
\end{equation}
로 정의한다. 전하 보존은
\begin{equation}
\mathcal{G}(V;q,T,\bm{\xi})=0
\label{eq:residual_zero}
\end{equation}
이라는 뜻이다.

여기서 내부 음극 전위 \(V_n\)를 처음 정의한다. \(V_n\)는 외부에서 미리 주어지는 함수가 아니라 식 \eqref{eq:residual_zero}을 만족하는 해이다.
\begin{equation}
V_n=\Phi(q,T,\bm{\xi}),
\qquad
\mathcal{G}(V_n;q,T,\bm{\xi})=0.
\label{eq:vn_root_definition}
\end{equation}
이를 풀어 쓰면
\begin{equation}
Q_{\cell}q
=
Q_{\bg}(V_n,T)
+\sum_{j=1}^{N_p}Q_{j,\tot}\xi_j
\label{eq:central_charge_balance_ch1}
\end{equation}
이다.

식 \eqref{eq:vn_root_definition}이 의미 있으려면 적어도 해가 존재해야 하고, 분석하는 branch에서 어떤 해를 따라갈지 정해야 한다. 가장 단순한 국소 조건은
\begin{equation}
\left|
\frac{\partial Q_{\bg}}{\partial V}(V_n,T)
\right|
\ge \epsilon_Q>0
\label{eq:bg_slope_nonzero}
\end{equation}
이다. 여기서 \(\epsilon_Q\)는 분석 구간에서 0이 아닌 기울기를 강제하기 위한 작은 양의 하한값이다. 이 조건은 \(Q_{\bg}\)가 전압 변화에 전혀 반응하지 않는 평평한 함수이면 작은 전하 불일치를 보정하기 위해 \(V_n\)가 무한히 크게 흔들릴 수 있다는 사실을 막는다. 여러 해가 존재하면 초기 조건에서 연속적으로 이어지는 branch를 택해야 한다.

이 단계에서 논리 순서는 분명하다. \(V_n\)가 먼저 있고 \(\xi_j\)를 만드는 것이 아니다. \(q\), \(T\), \(\bm{\xi}\)가 주어졌을 때 전하 보존을 만족시키는 전위가 \(V_n\)이다.

\section{평형 OCV는 특수한 해이다}

평형 상태에서는 각 transition 진행률이 같은 전위 \(V\)에서의 평형값을 따른다. 이 평형 진행률을
\begin{equation}
\xi_{j,\eqs}=\xi_{j,\eqs}(V,T)
\label{eq:xi_eq_general}
\end{equation}
로 둔다. 아직 \(V_n\)를 넣지 않고 \(V\)를 쓰는 이유는 정의 순서를 지키기 위해서이다.

평형 open-circuit 상태에서는 식 \eqref{eq:central_charge_balance_ch1}의 \(\xi_j\) 자리에 \(\xi_{j,\eqs}(V,T)\)를 넣는다.
\begin{equation}
Q_{\cell}q
=
Q_{\bg}(V,T)
+\sum_{j=1}^{N_p}Q_{j,\tot}\xi_{j,\eqs}(V,T).
\label{eq:equilibrium_balance}
\end{equation}
이 식을 \(V\)에 대해 푼 해를 평형 음극 OCV라 부른다.
\begin{equation}
V_{n,\OCV}(q,T)
=
\operatorname*{root}_{V}
\left[
Q_{\bg}(V,T)
+\sum_{j=1}^{N_p}Q_{j,\tot}\xi_{j,\eqs}(V,T)
-Q_{\cell}q
\right].
\label{eq:ocv_as_root}
\end{equation}

따라서 \(V_{n,\OCV}(q,T)\)는 이론의 출발점이 아니라 평형 조건을 추가했을 때 나오는 특수한 결과이다. 유한 전류, 저온, 이력 효과가 있으면 \(\xi_j\)가 \(\xi_{j,\eqs}\)를 즉시 따라가지 못하므로 같은 \(q\)에서도 \(V_n\neq V_{n,\OCV}\)가 될 수 있다.

\section{평형 transition 함수와 피크의 기본 형태}

상변이가 전압에 따라 한 상태에서 다른 상태로 넘어간다면, 평형 진행률은 전압에 대해 S자 형태를 갖는 것이 자연스럽다. 일반적인 설명을 위해 다음 형태를 예로 들 수 있다.
\begin{equation}
\xi_{j,\eqs}(V,T)
=
\left[
1+\exp\left(
-\frac{V-U_j(T)}{w_j(T)}
\right)
\right]^{-1}.
\label{eq:logistic_example}
\end{equation}
여기서 \(U_j(T)\)는 transition 중심 전위이고, \(w_j(T)>0\)는 transition이 전압축에서 얼마나 넓게 퍼지는지를 나타내는 폭이다. 이 식은 유일한 모델이라는 뜻이 아니다. ``전압이 중심 근처를 지날 때 진행률이 가장 빠르게 변한다''는 구조를 보이기 위한 예이다.

식 \eqref{eq:logistic_example}를 \(V\)로 미분하면
\begin{equation}
\frac{\partial \xi_{j,\eqs}}{\partial V}
=
\frac{1}{w_j(T)}
\xi_{j,\eqs}(1-\xi_{j,\eqs}).
\label{eq:logistic_derivative}
\end{equation}
이 값은 \(\xi_{j,\eqs}=1/2\), 즉 \(V=U_j(T)\) 근처에서 가장 크다. 따라서 평형 조건에서는 \(U_j(T)\) 부근에서 \(dQ/\dd V\) 피크가 생긴다.

더 열역학적인 표현으로 쓰면, 두 상태의 자유에너지 차이를 \(\Delta g_j(V,T)\)라고 둘 수 있다. 평형 진행률은
\begin{equation}
\xi_{j,\eqs}(V,T)
=
\left[
1+\exp\left(\frac{\Delta g_j(V,T)}{RT}\right)
\right]^{-1}
\label{eq:free_energy_occupancy}
\end{equation}
처럼 쓸 수 있다. 만약 \(\Delta g_j(V,T)\)가 전압에 대해 거의 선형이면 식 \eqref{eq:logistic_example}와 같은 꼴이 된다. 여기서 중요한 점은 온도가 평형 폭에도 영향을 준다는 것이다. 단순한 평형식만 보면 \(RT\)가 커질수록 transition은 전압축에서 넓어질 수 있다. 그러나 이것만으로는 저온에서 꼬리가 길어지는 현상을 설명하기 어렵다. 저온 tail은 뒤에서 보듯이 평형 폭보다 relaxation 속도 저하와 더 직접적으로 연결된다.

\section{비평형 진행률과 자기일관성}

유한 전류 또는 낮은 온도에서는 \(\xi_j\)가 평형값 \(\xi_{j,\eqs}(V_n,T)\)를 즉시 따라가지 못한다. 가장 단순한 상태 방정식은
\begin{equation}
\frac{\dd \xi_j}{\dd t}
=
k_j(V_n,T,q)
\left[
\xi_{j,\eqs}(V_n,T)-\xi_j
\right],
\qquad
k_j\ge 0.
\label{eq:xi_time_theory}
\end{equation}
이 식의 의미는 간단하다. 현재 진행률 \(\xi_j\)와 현재 전위에서 요구되는 평형 진행률 \(\xi_{j,\eqs}\) 사이에 차이가 있으면, 그 차이를 줄이는 방향으로 \(\xi_j\)가 움직인다. 그 속도는 \(k_j\)가 정한다.

\(\absI>0\)인 구간에서는 식 \eqref{eq:time_to_q}를 이용해
\begin{equation}
\frac{\dd \xi_j}{\dd q}
=
\frac{Q_{\cell}}{\absI}
k_j(V_n,T,q)
\left[
\xi_{j,\eqs}(V_n,T)-\xi_j
\right].
\label{eq:xi_q_theory}
\end{equation}
여기서 \(V_n\)는 독립 입력이 아니다. 식 \eqref{eq:vn_root_definition}에 의해
\begin{equation}
V_n(q)=\Phi(q,T(q),\bm{\xi}(q))
\label{eq:vn_phi_path}
\end{equation}
이다. 따라서 식 \eqref{eq:xi_q_theory}는 실제로
\begin{equation}
\frac{\dd \xi_j}{\dd q}
=
\frac{Q_{\cell}}{\absI}
k_j(\Phi(q,T,\bm{\xi}),T,q)
\left[
\xi_{j,\eqs}(\Phi(q,T,\bm{\xi}),T)-\xi_j
\right]
\label{eq:xi_q_closed}
\end{equation}
라는 자기일관 방정식이다.

이 구조에서 \(\bm{\xi}\)가 양쪽에 나타나는 것은 모순이 아니다. 미지수 \(x\)가 \(x=f(x)\)를 만족해야 한다는 방정식처럼, 미지의 함수 \(\bm{\xi}(q)\)가 자기 자신과 일관되는 경로를 찾아야 한다는 뜻이다. 적분형으로 쓰면 더 분명하다.
\begin{equation}
\xi_j(q)
=
\xi_j(q_0)
+
\int_{q_0}^{q}
\frac{Q_{\cell}}{\absI(s)}
k_j(V_n(s),T(s),s)
\left[
\xi_{j,\eqs}(V_n(s),T(s))-\xi_j(s)
\right]
\,\dd s.
\label{eq:xi_integral_theory}
\end{equation}
그리고 적분 안의 \(V_n(s)\)는 항상
\begin{equation}
Q_{\cell}s
=
Q_{\bg}(V_n(s),T(s))
+\sum_m Q_{m,\tot}\xi_m(s)
\label{eq:vn_inside_integral}
\end{equation}
를 만족한다. 이것이 되먹임 구조의 정확한 의미이다. 먼저 \(V_n(q)\) 곡선을 가정하는 것이 아니라, 전하 보존과 진행률 변화가 동시에 맞는 경로를 찾는 것이다.

\section{전하 보존식의 미분}

이제 ICA를 얻기 위해 식 \eqref{eq:central_charge_balance_ch1}를 미분한다. 등온이 아닌 경우를 먼저 쓰면
\begin{equation}
Q_{\cell}q
=
Q_{\bg}(V_n,T)
+\sum_jQ_{j,\tot}\xi_j.
\label{eq:balance_before_diff}
\end{equation}
양변을 \(q\)로 미분한다. 왼쪽은
\begin{equation}
\frac{\dd}{\dd q}(Q_{\cell}q)=Q_{\cell}
\label{eq:left_diff}
\end{equation}
이다. 오른쪽의 첫 항은 연쇄법칙으로
\begin{equation}
\frac{\dd}{\dd q}Q_{\bg}(V_n,T)
=
\frac{\partial Q_{\bg}}{\partial V_n}\frac{\dd V_n}{\dd q}
+
\frac{\partial Q_{\bg}}{\partial T}\frac{\dd T}{\dd q}
\label{eq:qbg_chain_rule}
\end{equation}
이다. transition 항은
\begin{equation}
\frac{\dd}{\dd q}
\left(
\sum_jQ_{j,\tot}\xi_j
\right)
=
\sum_jQ_{j,\tot}\frac{\dd \xi_j}{\dd q}
\label{eq:transition_diff}
\end{equation}
이다. 세 식을 합치면
\begin{equation}
Q_{\cell}
=
\frac{\partial Q_{\bg}}{\partial V_n}\frac{\dd V_n}{\dd q}
+
\frac{\partial Q_{\bg}}{\partial T}\frac{\dd T}{\dd q}
+
\sum_jQ_{j,\tot}\frac{\dd \xi_j}{\dd q}.
\label{eq:diff_balance_full}
\end{equation}
따라서
\begin{equation}
\frac{\dd V_n}{\dd q}
=
\frac{
Q_{\cell}
-
\dfrac{\partial Q_{\bg}}{\partial T}\dfrac{\dd T}{\dd q}
-
\sum_jQ_{j,\tot}\dfrac{\dd \xi_j}{\dd q}
}{
\dfrac{\partial Q_{\bg}}{\partial V_n}
}.
\label{eq:dvn_dq_general_theory}
\end{equation}

등온 구간에서는 \(\dd T/\dd q=0\)이므로
\begin{equation}
\frac{\dd V_n}{\dd q}
=
\frac{
Q_{\cell}
-
\sum_jQ_{j,\tot}\dfrac{\dd \xi_j}{\dd q}
}{
C_{\bg}
},
\qquad
C_{\bg}:=\frac{\partial Q_{\bg}}{\partial V_n}.
\label{eq:dvn_dq_iso_theory}
\end{equation}
여기서 \(C_{\bg}\)는 배경 저장량의 전압 미분이다. 선택한 전압 방향에 따라 부호 convention이 달라질 수 있으므로, 실제 branch에서는 \(C_{\bg}\)가 0에 가까워지지 않는다는 조건이 중요하다.

\section{평형 ICA 식}

평형에서는 \(\xi_j=\xi_{j,\eqs}(V_n,T)\)이다. 등온 평형 branch에서
\begin{equation}
\frac{\dd \xi_j}{\dd q}
=
\frac{\partial \xi_{j,\eqs}}{\partial V_n}
\frac{\dd V_n}{\dd q}
\label{eq:xi_eq_chain}
\end{equation}
이다. 이를 식 \eqref{eq:diff_balance_full}의 등온형에 넣으면
\begin{equation}
Q_{\cell}
=
C_{\bg}\frac{\dd V_n}{\dd q}
+
\sum_jQ_{j,\tot}
\frac{\partial \xi_{j,\eqs}}{\partial V_n}
\frac{\dd V_n}{\dd q}.
\label{eq:eq_diff_inserted}
\end{equation}
공통인 \(\dd V_n/\dd q\)를 묶으면
\begin{equation}
Q_{\cell}
=
\left[
C_{\bg}
+
\sum_jQ_{j,\tot}
\frac{\partial \xi_{j,\eqs}}{\partial V_n}
\right]
\frac{\dd V_n}{\dd q}.
\label{eq:eq_diff_grouped}
\end{equation}
따라서
\begin{equation}
\frac{\dd V_n}{\dd q}
=
\frac{Q_{\cell}}{
C_{\bg}
+
\sum_jQ_{j,\tot}
\dfrac{\partial \xi_{j,\eqs}}{\partial V_n}
}.
\label{eq:eq_dvn_dq}
\end{equation}

만약 관측 전압이 내부 전위와 같은 방향으로 \(V_{\obs}=V_n\)이라고 하면,
\begin{equation}
\frac{\dd Q_{\ext}}{\dd V_n}
=
\frac{\dd Q_{\ext}/\dd q}{\dd V_n/\dd q}
=
\frac{Q_{\cell}}{\dd V_n/\dd q}.
\label{eq:ica_from_reciprocal}
\end{equation}
식 \eqref{eq:eq_dvn_dq}를 넣으면
\begin{equation}
\boxed{
\frac{\dd Q_{\ext}}{\dd V_n}
=
C_{\bg}
+
\sum_jQ_{j,\tot}
\frac{\partial \xi_{j,\eqs}}{\partial V_n}.
}
\label{eq:eq_ica_decomposition}
\end{equation}

이 식이 평형 ICA의 핵심이다. 배경 항 \(C_{\bg}\) 위에 각 transition의 미분 항 \(Q_{j,\tot}\partial \xi_{j,\eqs}/\partial V_n\)이 더해진다. 즉 피크는 평형 진행률이 전압에 대해 빠르게 변하는 곳에서 생긴다.

\section{Apparent 전압과 실제 ICA}

실험에서 보이는 전압은 내부 음극 전위와 완전히 같지 않을 수 있다. 이를 가장 일반적으로
\begin{equation}
V_{\obs}
=
\Psi_{\obs}(V_n,q,T,\absI)
\label{eq:vobs_map}
\end{equation}
라고 쓰자. 여기서 \(V_{\obs}\)는 분석자가 ICA를 계산하는 전압 좌표이다. half-cell graphite 전위일 수도 있고, full-cell 전압에서 양극 기여와 다른 완만한 항을 분리하거나 포함한 유효 전압 좌표일 수도 있다. full-cell 전압을 그대로 사용할 때는 양극 전압 변화도 \(V_{\obs}\) 안에 함께 들어 있으므로, 흑연 피크 해석을 하려면 그 기여가 충분히 완만하거나 별도로 모델에 포함되어야 한다.

단순한 예는
\begin{equation}
V_{\obs}=V_n+s_I\absI R_n(q,T,\absI)
\label{eq:vobs_simple}
\end{equation}
이다. 여기서 \(s_I\)는 선택한 전압 방향 convention을 나타낸다.

식 \eqref{eq:vobs_map}를 \(q\)로 미분하면
\begin{equation}
\frac{\dd V_{\obs}}{\dd q}
=
\frac{\partial \Psi_{\obs}}{\partial V_n}
\frac{\dd V_n}{\dd q}
+
\frac{\partial \Psi_{\obs}}{\partial q}
+
\frac{\partial \Psi_{\obs}}{\partial T}
\frac{\dd T}{\dd q}
+
\frac{\partial \Psi_{\obs}}{\partial \absI}
\frac{\dd \absI}{\dd q}.
\label{eq:vobs_derivative_general}
\end{equation}
단순식 \eqref{eq:vobs_simple}에서 등온 정전류라면
\begin{equation}
\frac{\dd V_{\obs}}{\dd q}
=
\frac{\dd V_n}{\dd q}
+
s_I\absI\frac{\partial R_n}{\partial q}.
\label{eq:vobs_derivative_simple}
\end{equation}

측정 ICA는 항상
\begin{equation}
\boxed{
\frac{\dd Q_{\ext}}{\dd V_{\obs}}
=
\frac{Q_{\cell}}{\dd V_{\obs}/\dd q}
}
\label{eq:ica_observed}
\end{equation}
이고,
\begin{equation}
\boxed{
\frac{\dd V_{\obs}}{\dd Q_{\ext}}
=
\frac{1}{Q_{\cell}}
\frac{\dd V_{\obs}}{\dd q}.
}
\label{eq:dva_observed}
\end{equation}
이다. 따라서 \(dQ/\dd V\)를 이해하려면 먼저 \(dV_{\obs}/\dd q\)를 이해해야 하고, \(dV_{\obs}/\dd q\)의 중심에는 전하 보존식에서 나온 \(dV_n/\dd q\)가 있다. 이 관계는 분석 구간에서 \(V_{\obs}(q)\)가 국소적으로 미분 가능하고, \(\dd V_{\obs}/\dd q\)가 0으로 고정되지 않는 구간에서 해석한다.

\section{피크 면적은 왜 보존되는가}

어떤 transition \(j\)만 분리해서 보자. 그 transition이 전압 구간 \(V_a\le V_{\obs}\le V_b\) 안에서 만드는 전하 기여는
\begin{equation}
Q_j(V_{\obs})=Q_{j,\tot}\xi_j(V_{\obs})
\label{eq:transition_charge_vs_voltage}
\end{equation}
라고 쓸 수 있다. 그러면 그 transition의 ICA 기여는
\begin{equation}
\frac{\dd Q_j}{\dd V_{\obs}}
=
Q_{j,\tot}
\frac{\dd \xi_j}{\dd V_{\obs}}.
\label{eq:transition_ica_contribution}
\end{equation}
피크 면적은 이 값을 전압에 대해 적분한 것이다.
\begin{equation}
A_j
=
\int_{V_a}^{V_b}
Q_{j,\tot}
\frac{\dd \xi_j}{\dd V_{\obs}}
\,\dd V_{\obs}.
\label{eq:area_integral_start}
\end{equation}
미분과 적분이 서로 상쇄되므로
\begin{equation}
\boxed{
A_j
=
Q_{j,\tot}
\left[
\xi_j(V_b)-\xi_j(V_a)
\right].
}
\label{eq:area_integral_result}
\end{equation}
뒤의 괄호를
\begin{equation}
\Delta\xi_j=\xi_j(V_b)-\xi_j(V_a)
\label{eq:delta_xi_definition}
\end{equation}
라고 쓰면 \(A_j=Q_{j,\tot}\Delta\xi_j\)이다.

만약 해당 구간 안에서 \(V_{\obs}\)가 단조적으로 지나가고, transition이 거의 \(0\to 1\)로 완료된다면
\begin{equation}
A_j\approx Q_{j,\tot}.
\label{eq:area_complete_transition}
\end{equation}
따라서 같은 transition의 총 면적은 tail이 길든 짧든 거의 같다. tail이 길다는 것은 같은 면적이 더 넓은 전압 구간에 나누어져 있다는 뜻이지, 반드시 총량이 커졌다는 뜻이 아니다.

다만 이 결론에는 조건이 있다. 첫째, 배경 항이 적절히 분리되어야 한다. 둘째, 적분 구간이 transition 전체를 포함해야 한다. 셋째, 해당 구간에서 전압 좌표가 같은 값을 여러 번 되돌아가지 않는 단조 구간이어야 한다. 넷째, 여러 transition이 겹치면 각 \(A_j\)를 분리하기 어려울 수 있다. 이 조건이 깨지면 관측된 ``부분 면적''은 전체 transition 용량과 달라질 수 있다.

\section{피크 위치와 폭은 무엇을 뜻하는가}

식 \eqref{eq:transition_ica_contribution}에서 피크가 커지는 곳은 \(\dd \xi_j/\dd V_{\obs}\)가 큰 곳이다. 평형 조건에서 \(V_{\obs}\approx V_n\)이고 식 \eqref{eq:logistic_example}를 쓰면
\begin{equation}
\frac{\dd Q_j}{\dd V_n}
=
Q_{j,\tot}
\frac{1}{w_j}
\xi_{j,\eqs}(1-\xi_{j,\eqs}).
\label{eq:logistic_peak_contribution}
\end{equation}
이 함수는 \(V_n=U_j\) 부근에서 최대가 된다. 따라서 피크 위치는 transition의 중심 전위와 연결된다.

폭은 \(\xi_{j,\eqs}\)가 얼마나 넓은 전압 범위에서 변하는지와 연결된다. 식 \eqref{eq:logistic_derivative}에서는 \(w_j\)가 클수록 derivative peak가 넓어지고 낮아진다. 그러나 이것은 평형 폭이다. 실제 실험 피크의 폭과 tail은 여기에 비평형 지연이 추가되어 결정된다.

여러 transition이 가까운 전압에서 일어나면 ICA는
\begin{equation}
\frac{\dd Q_{\ext}}{\dd V_{\obs}}
\approx
\text{background}
+
\sum_j
Q_{j,\tot}
\frac{\dd \xi_j}{\dd V_{\obs}}
\label{eq:multi_peak_sum}
\end{equation}
처럼 보인다. 이때 서로 가까운 항들은 겹쳐서 하나의 넓은 피크처럼 보일 수 있다. 따라서 관측 피크 하나가 반드시 물리 transition 하나와 일대일 대응한다고 가정하면 안 된다.

\section{꼬리 형상은 왜 온도에 따라 달라지는가}

이제 사용자가 관측한 저온 long tail과 고온 short tail을 설명한다. 핵심은 \(\xi_j\)가 평형 목표값 \(\xi_{j,\eqs}\)를 얼마나 빨리 따라가는가이다.

식 \eqref{eq:xi_q_theory}에서 charge 좌표 기준 relaxation 계수를
\begin{equation}
\kappa_j(q,T)
=
\frac{Q_{\cell}}{\absI}
k_j(V_n,T,q)
\label{eq:kappa_definition}
\end{equation}
로 정의하자. 그러면
\begin{equation}
\frac{\dd \xi_j}{\dd q}
=
\kappa_j(q,T)
\left[
\xi_{j,\eqs}(V_n,T)-\xi_j
\right].
\label{eq:xi_q_kappa}
\end{equation}
이 식에서 \(\kappa_j\)가 크면 \(\xi_j\)는 작은 \(q\) 변화만으로도 평형값을 빠르게 따라간다. \(\kappa_j\)가 작으면 같은 전하가 더 많이 지나가도 \(\xi_j\)가 천천히 따라간다.

transition 뒤쪽 꼬리에서는 평형 목표값이 거의 최종값에 가까워졌다고 단순화할 수 있다. 방전 방향 완료값이 1이면
\begin{equation}
\xi_{j,\eqs}\approx 1.
\label{eq:tail_target}
\end{equation}
아직 완료되지 않은 잔여량을
\begin{equation}
u_j(q)=1-\xi_j(q)
\label{eq:unrelaxed_fraction}
\end{equation}
라고 두자. 또한 \(q_a\)를 꼬리 구간을 분석하기 시작하는 charge 좌표라고 두면, 식 \eqref{eq:xi_q_kappa}는
\begin{equation}
\frac{\dd u_j}{\dd q}
=
-\kappa_j(q,T)u_j
\label{eq:tail_decay_ode}
\end{equation}
가 된다. 이 식의 해는
\begin{equation}
u_j(q)
=
u_j(q_a)
\exp\left[
-\int_{q_a}^{q}\kappa_j(s,T(s))\,\dd s
\right].
\label{eq:tail_decay_solution}
\end{equation}
즉 tail은 지수적으로 사라지며, 사라지는 속도는 \(\kappa_j\)가 정한다.

이를 charge 좌표에서의 relaxation 길이로 쓰면
\begin{equation}
\ell_{q,j}(q,T)
=
\frac{1}{\kappa_j(q,T)}
=
\frac{\absI}{Q_{\cell}k_j(V_n,T,q)}.
\label{eq:relaxation_length}
\end{equation}
\(\ell_{q,j}\)는 무차원 좌표 \(q\) 위에서의 relaxation 길이이다. \(\ell_{q,j}\)가 크면 tail이 긴 \(q\) 범위에 걸쳐 남고, \(\ell_{q,j}\)가 작으면 tail이 짧게 끝난다.

이제 온도 의존성을 넣는다. 상변이가 활성화 자유에너지 장벽 \(\Delta G_{a,j}\)를 넘어 진행된다면, 단순한 형태로
\begin{equation}
k_j(V_n,T,q)
\sim
\nu_j(T)
\exp\left[
-\frac{\Delta G^\ddagger_j(V_n,T,q)}{RT}
\right]
\label{eq:activated_rate}
\end{equation}
라고 쓸 수 있다. 여기서 \(\nu_j(T)\)는 시도 빈도 또는 prefactor, \(\Delta G^\ddagger_j\)는 현재 전위와 상태에서의 활성화 자유에너지 장벽, \(R\)은 기체상수이다. 낮은 온도에서는 \(RT\)가 작아져 지수항이 작아지고, 보통 \(k_j\)가 감소한다. 그러면 식 \eqref{eq:relaxation_length}에 의해 \(\ell_{q,j}\)가 커진다. 결과적으로 같은 transition 용량이 더 긴 \(q\) 또는 전압 범위에 걸쳐 천천히 나타나므로 긴 tail이 생긴다.

반대로 높은 온도에서는 \(k_j\)가 커지고 \(\ell_{q,j}\)가 작아진다. 그러면 \(\xi_j\)가 평형 목표값을 빠르게 따라가므로 잔여량 \(u_j\)가 짧은 구간에서 사라지고 tail이 빨리 끝난다.

이 설명은 피크 면적 보존과 모순되지 않는다. 식 \eqref{eq:area_integral_result}에 따르면 총면적은 \(Q_{j,\tot}\Delta \xi_j\)이다. 온도는 \(\dd \xi_j/\dd V\)가 어느 전압 범위에 어떻게 분포되는지를 바꾸지만, 같은 구간에서 \(\Delta\xi_j\)가 같다면 면적은 같다.

\section{현재 전위 상태가 tail에 주는 영향}

온도만 중요한 것이 아니다. 현재 전위 상태 \(V_n\)도 중요하다. 식 \eqref{eq:xi_q_kappa}의 오른쪽에는
\begin{equation}
r_j(q)
=
\xi_{j,\eqs}(V_n(q),T(q))-\xi_j(q)
\label{eq:relaxation_residual}
\end{equation}
가 들어 있다. 이 값은 현재 전위에서 평형적으로 요구되는 진행률과 실제 진행률 사이의 차이이다.

만약 \(V_n\)가 transition 중심 전위 근처에 있으면 \(\xi_{j,\eqs}\)가 전압에 민감하게 변한다. 이때 \(\xi_j\)가 따라가지 못하면 \(r_j\)가 커지고, 그 차이가 ICA tail 또는 비대칭 broadening으로 나타난다. 반대로 \(V_n\)가 이미 transition에서 멀리 벗어났고 \(r_j\)도 거의 0이면 더 이상 큰 tail 신호가 남지 않는다.

따라서 tail 형상은 두 요소의 곱으로 이해할 수 있다.
\begin{equation}
\frac{\dd \xi_j}{\dd q}
=
\underbrace{\kappa_j(q,T)}_{\text{따라가는 속도}}
\underbrace{r_j(q)}_{\text{평형값과 실제값의 차이}}.
\label{eq:tail_two_factors}
\end{equation}
낮은 온도에서는 \(\kappa_j\)가 작아져 잔여량이 오래 남는다. 특정 전위 구간에서는 \(r_j\)가 커져 transition이 강하게 구동된다. 이 둘이 결합해 꼬리가 길어지거나 한쪽으로 늘어진 피크를 만든다.

\section{열역학과 동역학의 역할 분리}

여기서 주의할 점이 있다. 평형 열역학과 동역학 지연은 서로 다른 질문에 답한다.

\begin{longtable}{p{0.30\textwidth} p{0.62\textwidth}}
\toprule
항목 & 설명 \\
\midrule
\endhead
평형 열역학 & 현재 \(V_n,T\)에서 어떤 진행률 \(\xi_{j,\eqs}\)가 선호되는지 정한다. 즉 transition의 중심, 평형 폭, 가능한 저장량을 정한다. \\
동역학 지연 & 실제 \(\xi_j\)가 그 평형값을 얼마나 빨리 따라가는지 정한다. 즉 tail 길이와 비대칭 broadening을 정한다. \\
전하 보존식 & \(q\), \(V_n\), \(\bm{\xi}\)가 동시에 맞아야 하는 제약을 정한다. 따라서 \(\xi_j\) 지연은 곧 \(V_n\) 이동과 ICA 형상 변화로 연결된다. \\
\bottomrule
\end{longtable}

단순 평형식에서는 온도가 올라가면 \(\xi_{j,\eqs}\)의 전압 폭이 넓어질 수 있다. 그러나 사용자가 관측한 ``저온에서 긴 꼬리, 고온에서 짧은 꼬리''는 평형 폭 하나만으로 설명하기보다, 저온에서 \(k_j\)가 작아지고 \(\ell_{q,j}\)가 커지는 비평형 relaxation 효과로 설명하는 편이 더 논리적이다.

즉 이 문서의 해석은 다음과 같다.
\begin{equation}
\boxed{
\text{피크 위치와 평형 폭}
\leftarrow
\xi_{j,\eqs}(V,T),
\qquad
\text{tail 길이와 비대칭성}
\leftarrow
\ell_{q,j}=\frac{\absI}{Q_{\cell}k_j(V,T,q)}.
}
\label{eq:equilibrium_dynamic_split}
\end{equation}

\section{DVA와의 관계}

DVA는 ICA의 역수 구조를 갖는다. 식 \eqref{eq:ica_observed}와 \eqref{eq:dva_observed}에서
\begin{equation}
\frac{\dd Q_{\ext}}{\dd V_{\obs}}
\cdot
\frac{\dd V_{\obs}}{\dd Q_{\ext}}
=1
\label{eq:ica_dva_reciprocal}
\end{equation}
이다. 따라서 \(\dd V_{\obs}/\dd Q_{\ext}\neq 0\)인 미분 가능 구간에서 ICA의 큰 피크는 DVA의 작은 기울기 구간으로 나타난다. 다만 실제 데이터에서는 smoothing, noise, voltage orientation, full-cell 양극 기여 등이 있으므로 단순 역수만으로 모든 시각적 특징을 해석하면 안 된다. 이 장의 논리는 흑연 음극 기여를 분리해서 보는 이론적 기준이다.

\section{논리 점검표}

이 장의 논리 흐름은 다음 순서로 닫힌다.

\begin{enumerate}
\item 외부 전하 \(Q_{\ext}\)와 정규화 좌표 \(q\)를 먼저 정의한다.
\item ICA와 DVA를 측정 미분량으로 정의한다.
\item 흑연 저장량을 \(Q_{\bg}\)와 \(\sum_jQ_{j,\tot}\xi_j\)로 분해한다.
\item 전하 보존 residual \(\mathcal{G}=0\)을 쓴다.
\item 내부 전위 \(V_n\)를 residual의 해로 정의한다.
\item 평형 OCV는 \(\xi_j=\xi_{j,\eqs}\)를 대입한 특수한 해로 얻는다.
\item 비평형에서는 \(\xi_j\)가 \(\xi_{j,\eqs}\)를 시간 또는 charge 좌표에서 따라간다.
\item 이때 \(V_n\)는 매 순간 전하 보존식으로 연결되므로 자기일관 구조가 된다.
\item 전하 보존식을 미분하여 \(\dd V_n/\dd q\)를 얻는다.
\item 관측 전압 변환을 거쳐 \(\dd V_{\obs}/\dd q\)를 얻는다.
\item \(Q_{\ext}=Q_{\cell}q\)이므로 \(\dd Q_{\ext}/\dd V_{\obs}=Q_{\cell}/(\dd V_{\obs}/\dd q)\)이다.
\item transition 피크 면적은 \(Q_{j,\tot}\Delta\xi_j\)이다.
\item transition 피크 위치는 \(\dd\xi_j/\dd V\)가 가장 큰 곳과 연결된다.
\item 꼬리 길이는 \(\ell_{q,j}=\absI/(Q_{\cell}k_j)\)와 현재 residual \(r_j=\xi_{j,\eqs}-\xi_j\)로 설명된다.
\item 낮은 온도에서 \(k_j\)가 작아지면 \(\ell_{q,j}\)가 커져 tail이 길어진다.
\item 높은 온도에서 \(k_j\)가 커지면 \(\ell_{q,j}\)가 작아져 tail이 짧아진다.
\end{enumerate}

\section{현재 장의 한계}

이 장은 흑연 ICA/dQdV 피크 형상의 이론적 배경을 세우는 장이다. 따라서 실제 데이터에서 \(U_j\), \(Q_{j,\tot}\), \(w_j\), \(k_j\), \(R_n\) 등을 추정했다는 주장은 하지 않는다. 또한 특정 실험 데이터의 수치 오차나 최적 파라미터를 제시하지 않는다.

이 장에서 확정한 것은 논리 구조이다. 피크 면적은 transition 용량과 진행률 변화량에서 오고, tail 형상은 온도와 현재 전위 상태가 결정하는 relaxation 과정에서 온다. 후속 작업에서 데이터를 다루더라도 이 구조를 깨면 안 된다.

\end{document}
